\chapter{Methodology and Implementation}
\section{System Architecture}
The pipeline consists of raw-data validation, deterministic group assignment, split-manifest creation, development-only search, fixed-candidate refitting and holdout evaluation. Results are written as structured evidence rather than transcribed from console output. The same evaluator is used by both tuned methods, so preprocessing and fitness definitions remain consistent.

The main modules are \texttt{data.py} for ingestion and grouping, \texttt{evaluation.py} for fitting and scoring, \texttt{search.py} for Random Search, \texttt{de.py} for DE operators, and \texttt{stage3.py} for finalization and diagnostics. Verification logic audits saved records separately from optimization.
\section{Data Collection}
The dataset is the original UCI Concrete Compressive Strength spreadsheet and its readme \cite{uci}. Raw files and the downloaded archive are retained unchanged. The data are not scraped from a live service and are not augmented with another dataset. Source attribution includes both the UCI dataset DOI and Yeh's original publication. Table~\ref{tab:variables} lists the variables and units.

\begin{table}[htbp]\centering\caption{Dataset variables and units.}\label{tab:variables}\begin{tabularx}{\textwidth}{lX}\toprule Variable&Unit and role\\\midrule Cement; blast furnace slag; fly ash&kg/m$^3$, predictors\\Water; superplasticizer&kg/m$^3$, predictors\\Coarse aggregate; fine aggregate&kg/m$^3$, predictors\\Age&Days, predictor\\Compressive strength&MPa, target\\\bottomrule\end{tabularx}\end{table}
\section{Data Preprocessing}
The spreadsheet headers are mapped explicitly to canonical columns. The expected schema contains eight predictors and one target. Missing values are checked before splitting; no imputation is needed. All 1,030 rows are retained. There are 996 unique complete input tuples, so 34 rows repeat an existing eight-feature combination.

Rows with the same complete input tuple share a group. GroupShuffleSplit with seed 42 reserves 200 groups containing 204 rows for holdout; 796 groups containing 826 rows remain for development. Five GroupKFold splits are fixed inside development data. This isolates identical input tuples across fitting and evaluation boundaries while retaining their observed target values.
\section{Feature Engineering}
The model uses the eight original numerical inputs. It does not introduce technical indicators, polynomial interactions, feature selection or target transformations. Standardization is the only learned feature transformation. Keeping the feature representation fixed helps isolate the effect of hyperparameter optimization.
\section{SVR Model Architecture}
Each candidate instantiates a StandardScaler followed by an RBF SVR. The default comparator uses $C=1$, $\epsilon=0.1$ and gamma set to \texttt{scale}. Tuned candidates use decoded numerical values for all three parameters \cite{svr}. The scaler is refitted in every training fold and again on all development data for the final selected-model refit.
\section{Differential Evolution Implementation}
Population size is 15, mutation weight $F=0.8$, crossover rate $CR=0.9$ and the initialization-inclusive budget is 300. This permits one initial population and 19 complete trial generations. Every attempted fitness evaluation counts, including repeated positions. There is no hidden polishing stage or candidate cache that changes accounting.

\begin{algorithm}[htbp]\caption{Budget-controlled DE/rand/1/bin}\begin{algorithmic}[1]
\Require Bounds, population $N=15$, budget $B=300$, seed
\State Initialize $N$ candidates uniformly in log coordinates
\State Evaluate all candidates using five-fold development RMSE
\For{$g=1$ to $19$}
\State Freeze current population and fitness values
\For{each target $i$}
\State Select three distinct donors excluding $i$
\State Mutate with $F=0.8$ and clip to bounds
\State Cross over with $CR=0.9$, forcing one mutant coordinate
\State Evaluate trial; retain it if no worse than target
\EndFor
\State Commit the next population
\EndFor
\State \Return Candidate with lowest observed development RMSE
\end{algorithmic}\end{algorithm}
\section{Training Procedure}
For each seed, the optimizer receives the same saved five-fold development splits. The selected candidate minimizes mean fold RMSE. Each selected configuration is then refitted on all 826 development rows and evaluated on the same 204 holdout rows. Five seeds, 11, 22, 33, 44 and 55, are retained for each method.

The demonstration model is DE seed 11 because it has the lowest validation score under the predefined tie rule. Its holdout score is not used to choose it. Aggregate results include every predefined seed, rather than only this demonstration.
\section{Implementation Details}
The independent repository uses Python and \texttt{uv} with a committed lockfile. NumPy implements the explicit search operators; scikit-learn supplies scaling, SVR, splitting and metrics; pandas handles tabular data. Tests cover donor uniqueness, exclusion of the target, clipping, forced crossover, greedy selection, deterministic seeds and budget accounting.

Original final experiments were run on Windows with Python 3.13.1, NumPy 2.4.6, scikit-learn 1.6.1 and SciPy 1.15.3. Locked macOS verification used Python 3.12, NumPy 2.5.4, scikit-learn 1.9.1 and SciPy 1.18.1. These environments are distinguished in provenance; original results are preserved.
